ABSTRACT

\vspace{0.5cm}
\hrule
\vspace{0.5cm}

Digital audio workstations make detailed editing practical, but they store music as flat note events, so repetitive
structure is tedious to change and easy to lose.
Live-coding and scripting approaches recover programmability, yet they usually target real-time performance or
low-level event assembly rather than deterministic offline compilation to a portable score.

This thesis presents \textit{Motivo}, a domain-specific language and native compiler that keeps compositional structure
explicit in source and expands it into a Standard MIDI File without a language runtime at playback.
Motivo Studio wraps the toolchain in a browser-based editor for writing, compiling, and auditioning results.

We describe the language design, compiler implementation, studio application, and an empirical evaluation on synthetic
workloads.
The results support a bounded engineering claim: compact structural descriptions can be expanded into large timelines
within interactive compile times, without positioning Motivo as a replacement for established production tools.
